\section{A congruence-separated quartic norm in ten integers}
\label{sec:round9-integer}

The ten-integer route combines the new positive-gap, homogeneous-degree
source with the two-root congruence finalizer of the incoming
\emph{Congruence Norms} contribution. A further one-bit deletion removes
the outer leading block without losing the congruence required by the
arithmetic bridge. The result is an ordinary integer polynomial, not an
exponential Diophantine formula. All named registers below are abbreviations,
not additional unknowns. The exact coefficient payloads and literal circuit
are included in \texttt{integer/}.

\subsection{The leadless source and its calibration}

Use the gap source and its exact three-kernel certificate from the source
section. Write $W=\max_i w_i$, $H=\max_d(\deg D_d+dW)$, and
$A_T=\sum_jB^{q_j}$. Supply exactly
\[
 t,g_0,h_0,k_0,l,w,x,y,z,m\in\mathbb Z.
\]
In the primary mirror layout, row2 is additionally negated, an invertible
row operation; the kernel identity uses these oriented source rows.
Set $\kappa=2134940$, $\lambda=2^\kappa$, and define
{\small
\begin{align*}
 b&=6t+4,\quad g=bg_0,\quad h=bh_0,\quad k=b(3k_0+2),\quad B=\lambda b^2,\\
 \ell&=\sum_i B^{w_i},\quad u=\ell+g,\quad N_0=B^{W+1},\quad N=B^{W+H+2},\\
 G_c&=D_2(B)u^2+D_1(B)u+D_0(B),\quad M_c=N_0-1-(b-1)\ell,\\
 S_c&=g+N_0(2G_c+A_T),\quad T_c=M_c+N_0(B-2)A_T-2.
\end{align*}
}
The source ports for the bridge are
\begin{equation}\label{eq:r9-ten-source}
 \boxed{\quad X=N(S_c+T_c+1)-S_c-1,\qquad
 Y=2^{16\kappa-23}bB^{H+2W+9},\qquad
 \mathcal C=bB^W.\quad}
\end{equation}
No division appears in these expressions. The coefficient $2^{16\kappa-23}$
is materialized, not left as an unevaluated algorithmic parameter.

All signed assignments have $b\ne0$, $2\mid b$, and $b\equiv1\pmod3$.
Because all source weights are positive, direct reduction proves
\[
 X\text{ odd},\qquad b\mid X+1,\qquad b\mid Y,\qquad 3\mid X.
\]
Indeed $B,N,N_0\equiv1\pmod3$, $\ell\equiv31$, $A_T\equiv22$, so the
old high mask is $-22\equiv2$, the new $T_c$ is zero, and $X\equiv T_c$.

The deletion of $2$ from the mask is justified only after the finalizer has
recovered that $b$ is a positive even square. At that point $b\ge4$ and
$4\mid g$. The low base-$B$ digit of the old mask is $B-1$, while bit1 of
$S_c$ is zero. Thus clearing bit1 preserves binary disjointness and reduces
the mask population by exactly one. The usual source bounds give
$S_c+T_c<N$ for $0<g<2\mathcal C$.

If $b=2^s$ and $B=2^v$, then
\[
 \operatorname{pop}(T_c)=(W+23)v-31s-23.
\]
The two disjoint blocks of $X=(S_c+T_c)N+(N-1-S_c)$ give
\[
 \operatorname{pop}(X)=\log_2N+\operatorname{pop}(T_c)
 -\bigl(\operatorname{pop}(S_c)+\operatorname{pop}(T_c)
              -\operatorname{pop}(S_c+T_c)\bigr).
\]
The maximum is $(H+2W+25)v-31s-23=\log_2Y$. Consequently
\[
 Y\mid\binom{2X}{X}\quad\Longleftrightarrow\quad S_c\mathbin{\&}T_c=0,
\]
and the source certificate decodes a positive common source zero. This is
an exactly recalibrated threshold, smaller by four than the old leading-block
threshold. It does not assume $4\mid b$ on arbitrary signed inputs.

\subsection{The complete polynomial}

Define the division-free bridge registers
{\small
\begin{align*}
 L&=lY,&U&=2LX,&V&=4gwY,\\
 W_b&=bw,
 A&=U(V+1),&B_1&=2X+1,&C&=B_1+(A-2)h,\\
 P_{1/2}&=(LX)UV,&P&=2P_{1/2},&K&=X+1+k(P-2),\\
 S&=2A-5,&T&=3W_bC-2(W_b^2-1),&
 Q&=(P_{1/2}^2-1)K^2+1,\\
\end{align*}
}
{\small
\begin{align*}
 \Delta&=8\mathcal C^3gK^2
       -g^2\bigl(32(C-KL)^2\mathcal C^3+g^2K^2\bigr).
\end{align*}
}
The subscript $1/2$ names an explicitly multiplied polynomial; it is not an
instruction to divide an unknown. Put
{\small
\begin{align*}
 \xi&=zx,&D&=(A^2-4)C^2+4,&E&=C^2D\xi,\\
 F&=4(A^2-4)E^2+1,&G&=2(A+2)(A-2)^2E^2-1,\\
 \Omega&=16E,&H_c&=C+(B_1-C)F+\Omega Fy,&I&=(G^2-1)H_c^2+1,\\
 \Sigma&=\xi S^5,&\rho&=T+Sy,&
 \Theta&=(\rho^3-\rho S^2-S^3)(\rho^2+23S^2),\\
 r&=bQ,&a&=r+(3\Delta-1)(rz)^2,&d&=DFI,\\
 Z&=\Sigma m-\Theta.
\end{align*}
}
The final equation is
\begin{equation}\label{eq:r9-ten-final}
 \boxed{F_{10}=(Z^2-\Sigma^2(a+d))^2-4\Sigma^4ad=0.}
\end{equation}
The four fixed compiler ports are inherited as a whole from the original31
instance. The positive-gap transformations change its coordinates, not its
coefficient binding.

\subsection{Why the two-root norm is sufficient}

The normalizations imply $2b\mid U,A$, $b\mid K$, $3\mid K$, and
$Q\equiv1\pmod{b^2}$, $Q\equiv1\pmod3$. Also $2b\mid C+1$ and
$4b\mid D$. Therefore, without any positivity assumption,
\[
 a/b\text{ is odd},\qquad d/b\text{ is even},\qquad a\ne d.
\]
The odd integer $S$ is never zero. The quintic
$\pi(v)=(v^3-v-1)(v^2+23)$ has no rational root, so $\Theta\ne0$.
At a zero of \eqref{eq:r9-ten-final}, $\Sigma\ne0$. With
$q=Z/\Sigma$, the elementary identity
\[
 (q^2-a-d)^2=4ad
\]
and $a\ne d$ force $q\ne0$. The rational numbers
\[
 u_0=\frac{q^2+a-d}{2q},\qquad v_0=q-u_0
\]
satisfy $u_0^2=a$, $v_0^2=d$. They are integers because their squares
are integers. Hence both radicands are squares and $\Sigma\mid\Theta$;
the standard valuation argument gives $S\mid T$.

The coprime factors of
$a=r[1+(3\Delta-1)rz^2]$ are signed squares. Since $r\equiv1\pmod3$,
they must both be positive squares. The same argument for $r=bQ$, using
$\gcd(b,Q)=1$ and $b\equiv Q\equiv1\pmod3$, makes $b,Q$ positive
squares. Then $b\ge4$, and the second factor of $a$ forces $\Delta>0$.
The definition of $\Delta$ gives $K\ne0$ and $0<g<2\mathcal C$.
If $l=0$, then $P=0$ and $Q=1-K^2\le0$, a contradiction. Thus the
nonvanishing condition required by the inherited Sun bridge is recovered
rather than imposed as an additional polynomial factor. That bridge now
gives the central-binomial divisibility and hence a source zero.

For completeness, choose any sufficiently large $b=4^s$ in the source
encoding. Positive source weights give $b^2\mid g$ and
$X+1\equiv-g\pmod{b^2}$. The constructive bridge supplies the exact
Lucas values $C=u_{2X+1}(A)$ and $K=u_{X+1}(P)$. The recurrence and the
facts $X$ odd, $4b^2\mid P,A^2$ show that its existing witnesses satisfy
$b\mid h$, $b\mid k$, and $k/b\equiv2\pmod3$. Thus the displayed
restrictions lose no canonical solutions. Pell's theorem for the positive
nonsquare $(3\Delta-1)r$ chooses $z$ making $a$ square. The inherited
residue interpolation chooses $\xi$ divisible by $z$ and a suitable $y$
while keeping $d$ square and making $\Sigma\mid\Theta$. Then
$m=\Theta/\Sigma+\sqrt a+\sqrt d$ makes \eqref{eq:r9-ten-final} vanish.

These arguments compose with the source/refutation theorem to establish,
in ZFC, that the absence of integral solutions is equivalent to
$\operatorname{Con}(\mathrm{ZFC})$. The complete soundness and completeness
proof, including every source bound and congruence-class reconstruction,
is in \texttt{integer/PROOF.md}. The bridge references are Sun's
2021 Theorems4.12--4.13 and the modified certificate/interpolation results
in his 2026 ten-unknowns preprint. They remain explicit mathematical
dependencies, not claims of newly checked end-to-end kernel verification.

\subsection{Exact degree and evaluator accounting}

Write $r_X=4(W+H)+6$, $r_Y=2H+4W+19$, and $a_A=r_X+2r_Y+4$.
The dominant final term is uniquely $\Sigma^4d^2$, and
\[
 \deg d=88a_A+126,\qquad \deg\Sigma=5a_A+2.
\]
Thus the exact total degree is
\[
 \boxed{\deg F_{10}=196a_A+260=2352W+1568H+9668.}
\]
Full literal circuits, live-gate counts, actual-tuple modular replay, and
nonzero highest-degree coefficients over three primes independently check
the stated metrics. Global shared power chains reduce evaluation gates
without changing a single polynomial coefficient. The fixed threshold is
materialized as one leaf; charging the five constant-only gates in the older
evaluator convention adds five operations, with no effect on degree.
